\chapter{Blurry Vision}

\section*{Definition}
Blurry vision is reduced sharpness or clarity of sight. It may affect one or both eyes, be constant or intermittent, and occur at near or far distances. Sudden blurry vision can signal an emergency, whereas gradual blur commonly reflects refractive error or eye disease.

\section*{What Happens in Your Body}
Clear vision requires light to be focused by the cornea and lens onto a healthy retina, with signals transmitted through the optic nerve to the brain. Blur can result from an optical focusing problem, tear-film disturbance, corneal disease, cataract, retinal or optic-nerve disease, or disruption of visual pathways.

\section*{Causes}
\begin{itemize}
  \item Refractive error, dry eye, contact-lens problems, or eye strain
  \item Cataract, glaucoma, macular disease, diabetic retinopathy, or retinal detachment
  \item Corneal infection, inflammation, abrasion, or elevated eye pressure
  \item Migraine aura, stroke, transient ischemia, or optic neuritis
  \item Medication effects, high blood glucose, trauma, or hypertensive emergency
\end{itemize}

\section*{When to See a Doctor}
Seek emergency care for sudden vision change, a curtain or shadow, flashes with new floaters, severe eye pain, halos with nausea, or blurry vision with weakness, speech difficulty, facial droop, confusion, or severe headache. Prompt assessment is needed for pain, redness, injury, chemical exposure, pregnancy-related high blood pressure symptoms, or persistent unexplained blur.

\section*{Related Symptoms}
\begin{itemize}
  \item Eye pain, redness, discharge, light sensitivity, halos, flashes, or floaters
  \item Headache, dizziness, nausea, weakness, numbness, or speech difficulty
  \item Excessive thirst, frequent urination, or fluctuating vision with abnormal glucose
\end{itemize}

\section*{Self-Care / Home Management}
Do not drive if vision is unsafe. Remove contact lenses if irritating, use prescribed drops as directed, and avoid rubbing the eye. Flush chemical exposure with clean water continuously and obtain emergency care. Do not use steroid or numbing drops without an eye clinician.

\section*{How Your Doctor Will Evaluate You}
\begin{itemize}
  \item Assessment includes visual acuity with each eye, pupils, eye movements, visual fields, pressure, and slit-lamp or dilated retinal examination.
  \item Refraction, corneal testing, optical coherence tomography, or ocular ultrasound may identify the cause.
  \item Blood glucose, inflammatory markers, brain imaging, or vascular studies are selected for systemic or neurologic symptoms.
\end{itemize}

\section*{Treatment Options}
Treatment may include corrective lenses, lubrication, infection or inflammation treatment, pressure-lowering medicines, retinal procedures, cataract surgery, vascular emergency care, or management of diabetes and blood pressure. Treatment should be guided by the identified cause.

\section*{References}
\begin{thebibliography}{99}
\bibitem{blurry_vision:ref1} American Academy of Ophthalmology. Comprehensive Adult Medical Eye Evaluation Preferred Practice Pattern. AAO; 2024.
\bibitem{blurry_vision:ref2} American Academy of Ophthalmology. Diabetic Retinopathy Preferred Practice Pattern. AAO; 2024.
\bibitem{blurry_vision:ref3} Flaxel CJ, et al. Retinal vein occlusions Preferred Practice Pattern. \textit{Ophthalmology}. 2020;127:P288--P320.
\end{thebibliography}
